\begin{reviewinfo}
  
\begin{table}[tbh]
	\centering
	\begin{tabular}{|c|c|c|c|c|c|c|c|c|c|}
		\hline
		\hei{序号} &		\hei{评阅人} & 		\hei{职称}	& 	\hei{\makecell{导师\\类型}}	& 	\hei{工作单位}	
		& 	\hei{评分}	& 	\hei{结论}	& 	\hei{\makecell{答辩\\建议}} &		\hei{熟悉程度}& 	\hei{备注} \\

		\hline
		1 & 张三 & 教授& 博导& \makecell{XXX\\大学} &	95.8 &	达到	& A	& \makecell{有深入\\了解} &\makecell{行业\\专家}	  \\  \hline
	    2 & 李四 & 研究员& 硕导& \makecell{XXX\\大学} &	95 &	达到	& B	& \makecell{有深入\\了解} &	\\ \hline
	    3 & 王五 & 教授&  博导&\makecell{XXX\\大学} & 88& 达到& C & \makecell{有深入\\了解}& \\ \hline
	    4 & 赵六 & 教授&  博导& \makecell{XXX\\大学}& 88& 达到& C& \makecell{有深入\\了解}& \\ \hline
		\multirow{2}{*}{\makecell{\quad\\5}} & 孙六 & 教授&  博导&  \makecell{XXX\\大学} &	59 &	尚未达到	& D	& \makecell{有深入\\了解} & \\ \cline{2-10}
	    & 孙六 & 教授&  博导& \makecell{XXX\\大学} & 80 &	尚未达到	& B	& \makecell{有深入\\了解} & \makecell{复评\\结果} \\
		\hline
	\end{tabular}
\end{table}
\wuhao{
\noindent \textbf{说明}：

    1. 答辩建议选项包括4个：“A同意直接答辩”、“B同意修改后答辩”、“C须作较大修改后复评”、“D不同意答辩”。

    2. 熟悉程度选项包括3个：“有深入了解”、“比较熟悉”、“一般了解”。
\\

\color{red}{
\noindent 提醒（正式成文后删除）：

1. 评阅版论文删除此页。

2. 采用双盲评阅方式的学位申请人撰写的学位论文删除此页。

3. 评阅总分无需取整。

4. 工作单位填至学校、科研院所即可。

}


}



  
\end{reviewinfo}
% \emptydoublepage
